\section{Divisorial types and moving sections}\label{sec:divisorial-normalization-stability}

For a reference given by logarithms of holomorphic sections, removing its divisorial part identifies its ordinary singularity type with the bounded potentials of the residual background. The normalization scalar therefore depends only on the residual form, whereas convergence of the full potentials also requires control of the sections defining the divisors.

\Needspace{10\baselineskip}
In this section, $\theta$ is a smooth real $(1,1)$-form on $X$. Fix finitely many holomorphic line bundles $L_r$ on $X$ with smooth Hermitian metrics $h_r$, nonzero holomorphic sections $s_r$, and weights $t_r>0$. Put
\[
 \varphi=\sum_r t_r\log|s_r|_{h_r}^2,\qquad
 \Gamma=\sum_r t_r\operatorname{div}(s_r),\qquad
 \beta=\theta-\sum_r t_rc_1(L_r,h_r),
\]
where $c_1(L_r,h_r)$ is the Chern curvature form, and assume that the residual form $\beta$ is positive. The empty family is allowed, with $\varphi=0$ and $\beta=\theta$.

\begin{theorem}\label{thm:ordinary-divisor-normalization}
The ordinary singularity type of $\varphi$ is described by
\[
 \{u\in\PSH(X,\theta):u\sim\varphi\}
 =\{\varphi+w:w\in\PSH(X,\beta)\cap L^\infty(X)\},
\]
and for each such $w$,
\begin{equation}\label{eq:ordinary-divisor-measure}
 \mathrm{MA}^{\mathrm{ext}}_\theta(\varphi+w)=(\beta+\ddc w)^n.
\end{equation}
For every bounded measurable $F$, the equation \eqref{eq:relative-normalization} in this ordinary type has a unique constant and a solution unique modulo constants. Its normalization scalar $B_{\theta,[\varphi]}$ satisfies
\[
 B_{\theta,[\varphi]}(F)=B_\beta(F),
\]
so its dual functional and supporting sets coincide with $H_\beta$ and $\mathcal A_\beta$.
\end{theorem}
\begin{proof}
The Poincar\'e--Lelong formula gives $\theta+\ddc\varphi=[\Gamma]+\beta$, so every function on the right is $\theta$-plurisubharmonic and has the type of $\varphi$. Conversely, let $u\in\PSH(X,\theta)$ with $u\sim\varphi$. Then $u$ has the same generic Lelong numbers as $\varphi$ along the components of $\Gamma$, namely the coefficients of $\Gamma$, and \cref{lem:nonclosed-divisorial-removal} gives
\[
 \theta+\ddc u-[\Gamma]\geq0.
\]
The locally integrable difference $w=u-\varphi$ therefore has a $\beta$-plurisubharmonic representative. It is bounded off the divisor because $u\sim\varphi$, and the local submean property preserves both bounds across the divisor. Since $u$ and $\varphi+w$ agree almost everywhere, they are equal as quasi-plurisubharmonic functions.

Off the divisor, the currents $\theta+\ddc(\varphi+w)$ and $\beta+\ddc w$ agree, and locality identifies their Monge--Amp\`ere products there. The extended measure on the left gives the divisor zero mass, and so does the bounded product on the right, which is finite and charges no analytic set by \cref{sec:psh-bounded-products}. This proves \eqref{eq:ordinary-divisor-measure}. In the background $\beta$, the equation with bounded density has a continuous solution and a unique constant by \cref{thm:kn-continuous,prop:vanishing-exponent}. The density $\mathrm{e}^F$ is bounded above and below by positive constants, so \cref{cor:bounded-positive-density-uniqueness} gives uniqueness modulo constants. The correspondence transfers these assertions to the whole ordinary type and identifies the functionals.
\end{proof}

In particular, the comparison of backgrounds and the stability of supporting probabilities in \cref{thm:geometric-normalization-comparison,cor:geometric-adjoint-stability} apply whenever the positive residual forms satisfy the two-sided inequalities there; the location of the divisors does not enter these estimates. This statement concerns ordinary types. To interpret it in terms of model references, one also needs the envelope $P_\theta[\varphi]$ to differ from $\varphi$ by a bounded function. For a Hermitian $\theta$ and $n\geq2$, \cref{lem:one-sided-divisor-correspondence} below proves this bounded difference and the smaller-background model condition.

For a fixed positive residual $\beta$, uniformly bounded $F_j\to F$ in $L^1$ give uniform convergence of the residual solutions normalized by $\sup_Xv_j=0$. Indeed, continuity of the scalar gives uniformly bounded densities $\mathrm{e}^{B_\beta(F_j)+F_j}$ converging in every $L^p$, $p<\infty$. The compactness theorem \cref{thm:lp-uniform-compactness} gives uniform convergence of subsequences, continuity of bounded products identifies the limiting equation, and \cref{cor:bounded-positive-density-uniqueness} identifies every limit. Subtracting the means gives the same conclusion for solutions normalized to have mean zero.

\begin{theorem}\label{thm:moving-divisorial-potential-stability}\index{Divisor!moving sections}
Keep the line bundles and metrics fixed. Let nonzero sections $s_{jr}$ converge to $s_r\ne0$ in the finite-dimensional spaces of global sections, and let positive weights satisfy $t_{jr}\to t_r>0$. Given Hermitian forms $\beta_j$, define $\varphi_j$ and $\theta_j=\beta_j+\sum_rt_{jr}c_1(L_r,h_r)$ by the preceding formulas. Suppose either that $\beta_j\to\beta$ in $C^\infty$ for a Hermitian form $\beta>0$ and $F_j\to F$ in $C^\infty$, or that $\beta_j=\beta$ for all $j$ and the tests $F_j$ converge to $F$ in $L^1$ with a common $L^\infty$ bound. Then $\varphi_j\to\varphi$ in $L^1$, and the solutions $u_j\sim\varphi_j$ for $(\theta_j,F_j)$, normalized by $\sup_Xu_j=0$, converge in $L^1$ to the similarly normalized solution $u\sim\varphi$. For the residual solutions $v_j$ normalized to have mean zero, one also has
\[
 \sup_X(\varphi_j+v_j)\longrightarrow\sup_X(\varphi+v).
\]
The residuals converge smoothly in the first regime and uniformly in the second.
\end{theorem}
\begin{proof}
The convergence of the residual solutions follows from \cref{prop:smooth-background-potential-stability} in the first case and from the preceding compactness argument in the second. The fixed metrics and bounded weights give a common smooth lower bound for $\ddc\varphi_j$. The convergent sections give a common upper bound for $\varphi_j$ and convergence at a point outside all limiting divisors, so the suprema of the $\varphi_j$ are bounded above and below. By \cref{thm:normalized-compactness}, every subsequence has a further subsequence converging in $L^1$. Off the limiting divisors, convergence of the sections and weights gives local uniform convergence to $\varphi$; since those divisors have zero volume, every $L^1$ cluster point equals $\varphi$. Thus $\varphi_j\to\varphi$ in $L^1(X)$.

Put $C_j=\sup_X(\varphi_j+v_j)$ and $C=\sup_X(\varphi+v)$. The sums converge in $L^1$ and have a common smooth lower bound for their complex Hessians, since $\ddc v_j\geq-\beta_j$ and the residual forms are uniformly bounded. Applying \cref{lem:qpsh-hartogs} at maximum points gives $\limsup_jC_j\leq C$. A maximum point of $\varphi+v$ lies off the poles of $\varphi$, and evaluation there, using local uniform convergence, gives $\liminf_jC_j\geq C$.

It follows that $u_j=\varphi_j+v_j-C_j\to\varphi+v-C=u$ in $L^1$. For an empty family of divisors, the assertion reduces to the convergence of the residual solutions and of their suprema.
\end{proof}

\subsection{One-sided divisorial classes and sharp thresholds}
\label{subsec:divisorial-type-thresholds}

The ordinary-type correspondence extends to potentials with additional poles. It then identifies the largest subsolution coefficient in the whole one-sided class, even for merely integrable data. For smooth positive data, the prescribed divisor causes a strict loss relative to the global coefficient.

Keep the divisorial setting above, assume also that $\theta$ is a Hermitian form and $n\ge2$, and subtract a constant from $\varphi$ so that $\sup_X\varphi=0$. Write $S=\operatorname{Supp}\Gamma$ and $\psi=P_\theta[\varphi]$. For a nonempty class $\mathcal A\subset\PSH(X,\theta)$ and nonzero $0\le f\in L^1(\mathrm dV)$, define
\[
 b_{\theta,\mathcal A}(f)=\sup\{c\ge0:\text{some }u\in\mathcal A
          \text{ satisfies }\mathrm{MA}^{\mathrm{ext}}_\theta(u)
                                      \ge cf\,\mathrm dV\}.
\]
The three classes considered here are
\[
 \begin{aligned}
 \mathcal A_{\preceq}&=\{u\in\PSH(X,\theta):u\preceq\varphi\},\\
 \mathcal A_{\sim}&=\{u\in\PSH(X,\theta):u\sim\varphi\},\\
 \mathcal A_{\mathrm{env}}&=\mathcal E_{\mathrm{env}}(X,\theta,\psi).
 \end{aligned}
\]
The following lemma supplies the genuine reference $\psi$ used in the last definition.

\Needspace{12\baselineskip}
\begin{lemma}\label{lem:one-sided-divisor-correspondence}
The map $w\mapsto\varphi+w$ is a bijection from $\PSH(X,\beta)$ onto $\mathcal A_{\preceq}$, and
\[
 \mathrm{MA}^{\mathrm{ext}}_\theta(\varphi+w)
                         =\mathrm{MA}^{\mathrm{ext}}_\beta(w).
\]
Its bounded members correspond to $\mathcal A_{\sim}$. The envelope $\psi$ is genuine, satisfies $\varphi\le\psi\le\varphi+K$ for some finite $K$, meets the model condition \eqref{eq:als-model-condition}, and
\[
 \mathcal A_{\sim}\subset\mathcal A_{\mathrm{env}}
                         \subset\mathcal A_{\preceq}.
\]
\end{lemma}
\begin{proof}
If $u\preceq\varphi$, its generic Lelong numbers dominate the coefficients of $\Gamma$. By \cref{lem:nonclosed-divisorial-removal},
$\theta+\ddc u-[\Gamma]\ge0$. The locally integrable distribution $w=u-\varphi$ therefore has a canonical $\beta$-psh representative. Off $S$ it equals the difference and is bounded above. The identity $u=\varphi+w$ holds almost everywhere, hence everywhere for the canonical quasi-psh functions. Conversely, $w\in\PSH(X,\beta)$ has finite supremum, and
\[
 \theta+\ddc(\varphi+w)=\beta+\ddc w+[\Gamma]\ge0,
 \qquad \varphi+w\preceq\varphi.
\]
These maps are inverse. The bounded correspondence is \cref{thm:ordinary-divisor-normalization}. Off $S$ the currents agree; locality of the extended products identifies their measures there. Both measures give the analytic set $S$ zero mass by \cref{thm:arbitrary-extended-ma}, even when $w$ has additional poles or infinite mass. This proves the measure identity on $X$ without subtracting infinite pointwise values.

\Needspace{10\baselineskip}
For an envelope competitor $v\le0$ with $v\preceq\varphi$, the corresponding $w=v-\varphi$ satisfies $w\le-\varphi$ off $S$. Hence
\[
 \int_Xw^+\,\mathrm dV\le\int_X(-\varphi)\,\mathrm dV<\infty,
\]
uniformly in $v$. On a fixed finite cover by inner and outer coordinate balls, add smooth quadratic corrections making $w$ locally psh. Submean at a point attaining $\sup_Xw$, with uniform correction and volume-ratio bounds, gives
$\sup_Xw\le C_1\int_Xw^+\,\mathrm dV+C_2\le K$.
Thus $v\le\varphi+K$ canonically. The function $\varphi$ itself is a competitor, so the ordinary envelope construction gives a genuine $\theta$-psh function between $\varphi$ and $\varphi+K$. The upper bound survives upper regularization because $\varphi+K$ is upper semicontinuous. Ordinary equivalence gives $P_\theta[\psi]=P_\theta[\varphi]=\psi$. For $\varepsilon<1$ sufficiently close to one, $\beta-(1-\varepsilon)\theta>0$, so $\varphi$ is an $\varepsilon\theta$-psh competitor below $\psi$. Hence $P_{\varepsilon\theta}(\psi)$ is not identically $-\infty$, as required in \eqref{eq:als-model-condition}.

If $u\sim\varphi\sim\psi$, choose $K_u\ge0$ with $\psi-K_u\le u$. For each $t\ge1$, the $\theta$-psh function $a_t=\psi-tK_u$ obeys
\[
 t^{-1}a_t+(1-t^{-1})\psi=\psi-K_u\le u.
\]
This proves relative envelope-class membership in \cref{def:relative-envelope-class}, including the convention at $t=1$ and common poles. The other containment follows from its defining order $u\preceq\psi\sim\varphi$.
\end{proof}

\Needspace{12\baselineskip}
\begin{theorem}\label{thm:divisorial-type-thresholds}\index{Subsolution!prescribed divisorial threshold}
For every nonzero nonnegative $f\in L^1(\mathrm dV)$,
\[
 b_{\theta,\mathcal A_{\preceq}}(f)=b_\beta(f),
\]
and this coefficient is attained. If $f$ is smooth and strictly positive, then
\[
 b_{\theta,\mathcal A_{\sim}}(f)
 =b_{\theta,\mathcal A_{\mathrm{env}}}(f)
 =b_{\theta,\mathcal A_{\preceq}}(f)=c_\beta(f),
\]
with an exact finite equation attained in the ordinary type of $\varphi$. If also $\Gamma\ne0$, then
\[
 0<c_\beta(f)<c_\theta(f)=b_\theta(f).
\]
\end{theorem}
\begin{proof}
The correspondence lemma and \cref{thm:individual-density-subsolution} for the Hermitian form $\beta$ prove the first assertion, including attainment. This assertion allows additional poles and excess or infinite mass.
For smooth positive $f$, \cref{thm:tw-smooth} gives a smooth $w$ with $\beta+\ddc w>0$ and
$(\beta+\ddc w)^n=c_\beta(f)f\,\mathrm dV$.
The potential $\varphi+w$ belongs to the ordinary type and hence to all three classes. The class containments and the first assertion give equality of their coefficients; bounded-data normalization identifies $b_\beta(f)=c_\beta(f)$. The exact measure is finite.

It remains to prove strictness. Let $v$ be the smooth solution for $\theta$, and put
$\sigma=\theta+\ddc v$, $\eta=\beta+\ddc w$.
The one-sided subsolution gives $c_\beta(f)\le b_\theta(f)=c_\theta(f)$. If equality held, then on $X\setminus S$ the smooth function $a=\varphi+w-v$ would satisfy
\[
 \ddc a=\eta-\sigma,\qquad
 \ddc a\wedge Q=0,\qquad
 Q=\sum_{j=0}^{n-1}\eta^j\wedge\sigma^{n-1-j}.
\]
The form $Q$ is smooth and strictly positive, so it defines a real second-order operator $L$ with positive definite principal matrix and no zeroth-order term. Also $a\to-\infty$ near $S$, so its maximum is attained inside a component $U$ of $X\setminus S$. The strong maximum principle makes $a$ constant on $U$. To recall its local argument, put $z=\max_Ua-a\ge0$. If a positive region of $z$ had a boundary zero inside $U$, choose a ball in that region touching the zero set, with closure in a coordinate chart. On its outer half-annulus, the barrier
$e^{-A|x-y|^2}-e^{-AR^2}$ has strictly positive $L$ for large $A$. A sufficiently small multiple lies below $z$ on the two boundary spheres. A negative interior minimum of the difference would have nonnegative $L$, whereas its equation gives strictly negative $L$. Thus the comparison holds throughout the annulus and gives a positive inward derivative at the touching zero, contradicting $\nabla z=0$. Connectedness propagates the zero set to all of $U$.
The boundary of $U$ is nonempty and lies in $S$: otherwise $U$ would be both open and closed in connected $X$, although $S\ne\varnothing$. Constancy then contradicts $a\to-\infty$ at that boundary. This proves strictness; no semipositivity of the divisor's smooth curvature representative is required.
\end{proof}

\begin{corollary}\label{cor:critical-no-divisor}\index{Divisor!exclusion at the smooth threshold}
For smooth strictly positive $f$, no threshold subsolution has a nonzero divisorial component in $\theta+\ddc u$.
\end{corollary}
\begin{proof}
If such a component existed, retain a sufficiently small positive multiple $E$ of it. Choose a canonical logarithm $\ell_E$ and a smooth closed representative $\gamma_E$ with $[E]=\gamma_E+\ddc\ell_E$. The multiple can be chosen so that $\theta-\gamma_E>0$. Divisorial removal gives a canonical $(\theta-\gamma_E)$-psh representative of $u-\ell_E$, whose finite supremum implies $u\preceq\ell_E$. The preceding theorem would then give
\[
 b_\theta(f)\le c_{\theta-\gamma_E}(f)
                         <c_\theta(f)=b_\theta(f),
\]
a contradiction.
\end{proof}
The later rigidity result \cref{cor:smooth-dual-exact-equation} strengthens this obstruction by determining every smooth-data threshold subsolution through the maximizing Gauduchon form.

\Needspace{15\baselineskip}
\begin{example}\label{ex:projective-divisor-threshold-loss}
Let $(X_0,h_0)$ be a connected compact Hermitian $m$-fold, let $k\ge2$, and let $\kappa$ be the mass-one Fubini--Study form on $\mathbb P^k$. On $X_0\times\mathbb P^k$, choose a smooth positive volume form $\mathrm dV$. With pullbacks implicit, put
\[
 \begin{aligned}
 N&=m+k,& C&=\binom Nm,& V_0&=\int_{X_0}h_0^m,\\
 H&=h_0+\kappa,& \mu&=(h_0^m/V_0)\wedge\kappa^k,& R&=\mu/\mathrm dV.
 \end{aligned}
\]
For $D=\{Z_0=0\}$ pulled back to the product, take $E=\lambda D$, $0<\lambda<1$, with
\[
 \ell_E=\lambda\log\frac{|Z_0|^2}{\sum_{i=0}^k|Z_i|^2},
 \qquad H_E=h_0+(1-\lambda)\kappa.
\]
The global threshold of $R$ is $CV_0$, while its ordinary, relative envelope and one-sided thresholds prescribed by $E$ are all $CV_0(1-\lambda)^k$, attained by $\ell_E$. If $v_-(h_0)=0$, both $H$ and $H_E$ are nonclosed and have zero lower mass.
\end{example}
\begin{proof}
The Poincar\'e--Lelong formula gives $[D]=\kappa+\ddc(\ell_E/\lambda)$. Dimensions and the mass-one projective factor give
\[
 H^N=CV_0R\,\mathrm dV,\qquad
 H_E^N=CV_0(1-\lambda)^kR\,\mathrm dV.
\]
Thus the zero residual solves each smooth positive equation and fixes its unique scalar. The theorem identifies all three restricted coefficients, and the correspondence gives
$\mathrm{MA}^{\mathrm{ext}}_H(\ell_E)=H_E^N$.
Under the factor hypothesis, pullbacks of bounded $h_0$-psh potentials with mass tending to zero have masses multiplied respectively by $C$ and $C(1-\lambda)^k$. Bounded submersion pullback and the product expansion therefore give zero lower mass for both backgrounds. A closed positive $h_0$ would instead have every bounded mass equal to $V_0$ by Stokes, so both backgrounds are nonclosed.
\end{proof}

